\chapter{Method}
\section{Markov chain Monte Carlo}
If we want to sample directly from the input space, 
we must somehow determine which samples to sample. In Markov chain Monte Carlo, 
we start with an initial sample $X_1$. The marginal distribution is called the 
``initial distribution''~\cite{geyer2011introduction}.
Each subsequent sample is drawn based on the previous one using a transition probability distribution. Crucially, in a Markov chain, the next state $X_{n+1}$ depends only on the current state $X_n$, rather then on the full history of the chain.
Since we do not want random samples, we require that the chain convergeing to a target distribution $\pi$.
To make it more concrete, we have to understand what it means for a distribution to remain unchanged
as the chain evolves, this leads us to the concept of stationary.
Stationary means that the conditional distribution for a 
subset of the chain does not depend on where the subset is taken from.
We can be sure that a chain is stationary by checking if it is reversible. 
For a reversible chain, the distribution for $X_{n+1},X_{x+2},\dots,X_{x+k}$ and $X_{n+k},\dots,X_{n+2},X_{n+1}$ are the same.
Often a burn-in is used, i.e. 
the first $n$ values of the chain are discarded with the aim of throwing away these samples before the chain converges. 
According to Geyer ~\cite{geyer2011introduction}, 
burn-in is unnecessary and starting from a reasonable point is just as effective. 
Therefore, in this paper, it is omitted in all examples.
As said previously, we select the next sample based on the transition probabilty distribution, but we have not yet defined how exactly it is defined. In this work we will keep our focus on the Metropolis-Hastings algorithm.

\subsection{Metropolis-Hastings algorithm}
The Metropolis-Hastings algorithm is one possible implementation of constructing the Makrov chain in MCMC and 
is developed by Hastings~\cite{hastings1970}.
Let $\pi(x)$ be the density we want to take samples from, called the target density.
Then, the algorithm requires a distribution proportional to the target denisity, $\pi(x)\propto\tilde{\pi}(x)$.
To determine the next state in the Makrov chain, it also requires a proposal kernal $q$.

The transition from the Makrov state $X_t$ to $X_{t+1}$ is then defined as follows.
Sample a new State  $X_{t+1}$ from $q(y | X_t)$. 
Accept or reject the State with the acceptance probability
\[ 
    \alpha(X_t, X_{t+1}) = 
    \begin{cases}
        \min \left\{ \frac{\tilde{\pi}(X_{t+1}) q(X_t|X_{t+1} )}{\tilde{\pi}(X_t) q(X_{t+1}|X_t)}, 1 \right\} & \text{if } \tilde{\pi}(X_t)q(X_{t+1}|X_t)>0 \\
        1 & \text{otherwise}
    \end{cases}     
\]
If accepted, then set $X_{t+1}=X_{t+1}$, otherwise $X_{t+1}  = X_t$ and restart the process. 
Importantly this also means if the sample is rejected, the sample before is twice in the chain.

Metropolis-Hastings is a further developlment of the Metropolis algorithm, which was introduced by Metropolis et al.~\cite{metropolis1953} and required $q$ to be symmetrical.

In Random Walk Metropolis, the proposals are drawn with i.i.d and fixed symmetric densities $Z_i$,
for example $Z_i\sim N(0,\sigma^2 I_d)$.
Each new proposal therefore has the form $X_{i+1} = X_i + Z_{i+1}$~\cite{rosenthal2011optimal}.
\subsection{Hamiltonian Dynamics}
\subsection{Hamiltonian Monte Carlo}
Hamiltonian Monte Carlo is an adaptation of Metropolis-Hastings, which uses 
Hamiltonian dynamics to sample from a unknown distribution, in bayesian conext the posterior distribution.
It was originally published under the name of Hybrid Monte Carlo by Duane et al.~\cite{DUANE1987216}, but 
it became known under the name of Hamiltonian Monte Carlo, due to the more specific formulation.
Since Hamiltonian dynamics normally occurs in the physical context, the terms are also adopted in the machine learning context. 

As we saw in the previous section , Hamiltonian dynamics requires one $d$-dimensional vector, describing the position $q_i$, with $i=1,...,d$ entries. 
The vector here is the one we want to sample from, and consists in difference to normal Hamiltonian dynamics out of the network parameters. 
The probability denisity for $q$ is defined under the canonical distribution with:
\begin{align}
P(q)\propto \exp(-U(q))\label{eq:2.1}
\end{align}


for some energy function $U(q)$.
Since we want to use dynamical methods, we require
the "momentum" vector, $p_i$, with $i=0,...,d$ entries. 
Together, the distribution combines to:
\begin{align}
P(q,p)\propto \exp(-H(q,p))\label{eq:2.2}
\end{align}
Where $H(q,p) = U(q)+K(p)$ is called the \textit{Hamiltonian}, 
$U(q)$ is the potential energy, and $K(p)$ is the kinetic energy. 
If we insert this into the equation \ref{eq:2.1}, we get:
\begin{align}
    P(q,p)\propto\exp(-U(q))\exp(-K(p))\label{eq:2.3}
\end{align}

Since we are in a bayesian context, and the main focus is to sample from the posterior distribution, we can display $U(q)$ as: 
\[
U(q) = -\log[\pi(q)L(q|D)]
\]
with $\pi(q)$ is the prior density, and $L(q|D)$ is the likelihood.
\todo{warum darf ich U/K so definieren? -> Formel K: $\sum_{i=1}^d \frac{p_i^2}{2m_i}$ \& Neal hier zitieren!}



% Similarly, we require
% the "momentum" vector, $p_i$, with $i=0,...,d$ entries. Obviously the momentum is not the actual movement of particals here (especially since there are no particals in this use-case),
% but it is defined as the gradient of some function $U(q)$, called the potential energy. 
\begin{align*}
    \frac{dq_i}{d_t}&=\frac{\partial H}{\partial p_i} \\
    \frac{dp_i}{d_t}&=-\frac{\partial H}{\partial q_i}
\end{align*}

for $i=1,...,d$~\cite{MR2858447}. 


  

\begin{itemize}
    \item p two dimensional vektor momentum -> $m\cdot v$
    \item q two dimensional vektor position 
    \item $U(q)$ potential energy $U(q)\propto$ height of surface and kinetic energy
    \item kinetic energy $K(p)=\frac{\abs{p}^2}{2m}$
    \item constant velocity $v=\frac{p}{m}$
\end{itemize}   
\subsection{Metropolis adjusted Langevin algorithm}
The Metropolis adjusted Langevin algorithm (MALA) was first proposed by 
Besag~\cite{besagComment} and further evaluated by Roberts and Tweedie~\cite{roberts1996exponential}.

